%% 5-3 かぎりなく近づく：極限と接線（第11回：ページ制限なし。偶数ページ）
\bkchapter{5}{3}{かぎりなく近づく：\\極限と接線}{\bklevel{3}}{かぎりなく近づく：極限と接線}
\bkmichishirube{3-5,4-2,5-2}{5-4,5-8}

\begin{goal}
\begin{itemize}[label={\color{bkTheme}●}, leftmargin=1.6\zw, itemsep=0pt, topsep=1pt]
\item \bklevel{3}\enspace $\lim$ の記号を読み、極限の値を求められる。
\item \bklevel{3}\enspace 微分係数を、接線の傾きとして求められる。
\item \bklevel{3}\enspace 導関数を求め、接線の式を書ける。\bkidx{せっせんのしき}{接線の式}
\end{itemize}
\end{goal}

\begin{hatena}
\noindent\begin{minipage}[c]{0.55\linewidth}
スケートボードのランプ（すべり台のような曲がった坂）を考えます。坂の急さは、場所ごとにちがいます。下のほうはゆるやかで、上へ行くほど急です。

では、点Pのところの「坂の急さ」は、どんな数で表せばよいでしょう。まっすぐな坂なら傾きで表せました（3-5）。でも、曲がった坂には、決まった傾きがありません。
\end{minipage}\hfill
\begin{minipage}[c]{0.42\linewidth}
\begin{tikzpicture}
\begin{axis}[width=58mm, height=46mm, axis lines=none, xmin=-0.2, xmax=2.3, ymin=-0.3, ymax=4.6]
  \fill[bkYellow!25] (axis cs:0,0) -- plot[domain=0:2, samples=40] ({\x},{\x^2}) -- (axis cs:2,0) -- cycle;
  \addplot[bkcurve, domain=0:2] {x^2};
  \draw[bkGray, line width=1pt] (axis cs:-0.2,0) -- (axis cs:2.3,0);
  \addplot[only marks, mark=*, mark size=2pt, bkInk] coordinates {(1.2,1.44)};
  \node[anchor=south east, font=\sffamily\small] at (axis cs:1.2,1.44) {P};
  \addplot[bkcurve2, line width=1pt, dashed, domain=0.65:1.72] {1.44+2.4*(x-1.2)};
  \node[anchor=west, font=\sffamily\scriptsize, text=bkGray] at (axis cs:1.78,2.9) {?};
\end{axis}
\end{tikzpicture}
\end{minipage}
\end{hatena}

5-2では、時間のあいだ $h$ を0に近づけて、瞬間の速さを求めました。同じ考え方で、曲がったグラフの1点での「傾き」が決められます。この章では、その考え方に名前と記号をつけます。

\clearpage
%% ---------------- やってみよう・思い出そう ----------------
\begin{yattemiyou}[グラフにぴったりふれる直線を引こう]
\textbf{用意するもの}\par
方眼紙、定規
\begin{enumerate}[label={\bkcircled[bkGreen]{\arabic*}}, leftmargin=2.4\zw, itemsep=1pt, topsep=3pt]
\item 下の表の点を方眼紙にとり、なめらかな線でつないで、$y=x^2$ のグラフをかこう。1目もりは0.5にする。
\item 点P（1，1）と点Q（2，4）を結ぶ直線を引き、傾きを求めよう。
\item Qを（1.5，2.25）にかえて、同じように傾きを求めよう。
\item 定規を点Pに当て、グラフにぴったりふれるように直線を引こう。その傾きは、いくつになったかな？
\end{enumerate}
\begin{bkcenter}
\begin{bktable}{colspec={Q[l,m,wd=14mm] *{9}{Q[c,m,wd=9mm]}}, rowsep=1pt, colsep=3pt}
$x$ & $-2$ & $-1.5$ & $-1$ & $-0.5$ & 0 & 0.5 & 1 & 1.5 & 2 \\
$y=x^2$ & 4 & 2.25 & 1 & 0.25 & 0 & 0.25 & 1 & 2.25 & 4 \\
\end{bktable}
\end{bkcenter}
\end{yattemiyou}

\noindent 傾きは「たての変化量÷横の変化量」です（3-5）。②の傾きは $\dfrac{4-1}{2-1}=\dfrac{3}{1}=3$、③は $\dfrac{2.25-1}{1.5-1}=\dfrac{1.25}{0.5}=2.5$ です。QをPに近づけると、傾きは小さくなっていきます。④で引いた直線の傾きは、2に近くなったはずです。この章で、ちょうど2になることを確かめます。

\begin{omoidasou}
\bkrecall{5-2}{$x=5t^2$ の、1秒から $1+h$ 秒までの平均の速さは $10+5h$。$h$ を0に近づけると10。}
\bkrecall{3-3}{$f(x)$ は「関数 $f$ に $x$ を入れた値」。\\$f(x)=x^2$ なら $f(3)=3^2=9$、$f(a+h)=(a+h)^2$。}
\bkrecall{3-5}{傾きが $m$ で、点 $(x_1,\ y_1)$ を通る直線の式は $y-y_1=m(x-x_1)$。}
\bkrecall{4-2}{円の接線は、円に1点だけでふれる直線。接線は、ふれた点で半径に垂直。}
\end{omoidasou}

\noindent 3-3の $f(x)$ は、「$x$ を入れると、きまりにしたがって数が出てくる箱」のようなものです。$f(x)=x^2$ という箱に3を入れると9が、$a+h$ を入れると $(a+h)^2$ が出てきます。5-1・5-2で使った $x(1.1)$ も、同じ書き方でした。

\begin{bkfigure}
\begin{tikzpicture}[x=1mm, y=1mm]
  \node[bk node, minimum width=30mm, minimum height=12mm, font=\sffamily\small, fill=bkTheme!10] (f) at (60,0) {$f$\enspace（2乗する）};
  \node[font=\sffamily\small] (in1) at (14,5) {3};
  \node[font=\sffamily\small] (in2) at (14,-5) {$a+h$};
  \node[font=\sffamily\small] (o1) at (108,5) {$f(3)=9$};
  \node[font=\sffamily\small] (o2) at (108,-5) {$f(a+h)=(a+h)^2$};
  \draw[bk link] (in1) -- ([yshift=5mm]f.west);
  \draw[bk link] (in2) -- ([yshift=-5mm]f.west);
  \draw[bk link] ([yshift=5mm]f.east) -- (o1);
  \draw[bk link] ([yshift=-5mm]f.east) -- (o2);
\end{tikzpicture}
\bkfigcaption{関数 $f(x)=x^2$ は、入れた数を2乗して出す箱。}
\end{bkfigure}

%% ---------------- 数学の道具 ----------------
\bksection{doguu}{数学の道具}

\bksubsection{1}{かぎりなく近づく：極限と $\lim$}

\leavevmode\label{bk:5-3:lim}5-2では、「$h$ を0に近づけると、$10+5h$ は10に近づく」と考えました。表にすると、次のようになります。

\begin{bkcenter}
\begin{bktable}{colspec={Q[l,m,wd=22mm] *{5}{Q[c,m,wd=17mm]}}, rowsep=1.5pt}
$h$ & 0.1 & 0.01 & 0.001 & 0.0001 & …… \\
$5h$ & 0.5 & 0.05 & 0.005 & 0.0005 & …… \\
$10+5h$ & 10.5 & 10.05 & 10.005 & 10.0005 & → 10 \\
\end{bktable}
\bktabcaption{$h$ を0に近づけると、$10+5h$ は10に近づく}
\end{bkcenter}

$h$ は0にはなりません。0でない数のまま、0にいくらでも近くなります。このようすを「\textbf{$h$ を0にかぎりなく近づける}」といい、そのとき式の値が近づいていく先の10を、\textbf{極限}（きょくげん）といいます。記号で、次のように書きます。\bkidx{きょくげん}{極限}
\[ \lim_{h\to0}(10+5h)=10 \]

\begin{bkfigure}
\begin{tikzpicture}[x=1mm, y=1mm]
  % lim の書き方。式を3つの部分に分けて並べ、説明は各部分の真上・真下に置く（線は交差しない）
  \node[font=\Large, inner sep=1pt] (A) at (0,0) {$\displaystyle\lim_{h\to0}$};
  \node[font=\Large, anchor=base west, inner sep=1pt] (B) at ([xshift=2mm]A.base east) {$(10+5h)$};
  \node[font=\Large, anchor=base west, inner sep=1pt] (C) at ([xshift=2mm]B.base east) {$=\ 10$};
  \node[anchor=south east, font=\sffamily\footnotesize, text=bkBlue] (a1) at ([yshift=7mm, xshift=3mm]A.north) {リミット（かぎりなく近づけると）};
  \node[anchor=north, font=\sffamily\footnotesize, text=bkBlue] (a2) at ([yshift=-7mm]A.south) {$h$ を0に近づけるとき};
  \node[anchor=south west, font=\sffamily\footnotesize, text=bkGreen!70!black] (b1) at ([yshift=7mm, xshift=-3mm]B.north) {この式の値が近づく先は};
  \node[anchor=north, font=\sffamily\footnotesize, text=bkOrange!80!black] (c1) at ([yshift=-7mm, xshift=6mm]C.south) {10（極限）};
  \draw[bkBlue, line width=0.5pt] ([yshift=5.5mm]A.north) -- ([yshift=1.5mm]A.north);
  \draw[bkBlue, line width=0.5pt] ([yshift=1.5mm]a2.north) -- ([yshift=-1.5mm]A.south);
  \draw[bkGreen!70!black, line width=0.5pt] ([yshift=5.5mm]B.north) -- ([yshift=1.5mm]B.north);
  \draw[bkOrange!80!black, line width=0.5pt] ([yshift=1.5mm, xshift=-6mm]c1.north) -- ([yshift=-1.5mm]C.south);
\end{tikzpicture}
\bkfigcaption{$\lim$ の式の読み方}
\end{bkfigure}

「リミット、$h$ が0に近づくとき、$10+5h$ は10」と読みます。$\lim$ は、英語で「限り・限界」という意味の limit の頭の3文字です。$\lim$ の下の $h\to0$ は「$h$ を0に近づける」という意味で、矢印「→」は「近づく」と読みます。\bkidx{りみっと}{$\lim$（リミット）}

\noindent\textbf{「$=10$」と書いてよいの？}\par
$10+5h$ は、$h$ が0でない限り、ちょうど10にはなりません。でも、$\lim$ の式は「$10+5h$ の値」ではなく、「近づく先」を表しています。近づく先は、ちょうど10です。だから「$=10$」と書いてよいのです。

\begin{seikaku}
\textbf{「かぎりなく近づく」とは}\par
$h$ を0でない数のまま小さくしていくと、$10+5h$ と10の差 $5h$ は、0.001より小さく、0.000001より小さく……と、どんなに小さい数を決めても、それより小さくできます。このことを「$10+5h$ は10にかぎりなく近づく」といいます。大学の数学では、この言い方をもっと正確な式で書きます。
\end{seikaku}

\noindent\textbf{分数の形の極限}\par
$\displaystyle\lim_{h\to0}\frac{6h+h^2}{h}$ のように分母に $h$ があるときは、$h$ に0を入れると $\dfrac{0}{0}$ になって計算できません。そこで、$h$ を0に近づける前に、分子の1つずつを $h$ でわります（$h$ は0ではないので、わってよい。2-12）。\bkidx{ぶんぼ}{分母}
\[ \lim_{h\to0}\frac{6h+h^2}{h}=\lim_{h\to0}\Bigl(\frac{6h}{h}+\frac{h^2}{h}\Bigr)=\lim_{h\to0}(6+h)=6 \]
数で確かめると、$h=0.1$ なら $\dfrac{0.6+0.01}{0.1}=\dfrac{0.61}{0.1}=6.1$、$h=0.01$ なら $\dfrac{0.06+0.0001}{0.01}=6.01$ で、6に近づいています。

\bkkigocards{\bkkigocard{$\lim$}{リミット}{かぎりなく近づけたときの、近づく先の値（極限）}{limit}{5-3}}
{\bkkigocard{$\dfrac{\dd y}{\dd x}$}{ディーワイ・ディーエックス}{$y$ を $x$ で微分したもの（導関数）}{derivative}{5-3}}

\bksubsection{2}{曲線の接線}

\leavevmode\label{bk:5-3:sessen}円の接線は、円に1点だけでふれる直線でした（4-2）。曲がったグラフにも、1点でぴったりふれる直線を考えます。

グラフの上に点Pをとり、もう1つの点QをPの近くにとって、直線PQを引きます。Qをグラフにそって、Pにかぎりなく近づけていくと、直線PQはある1本の直線に近づいていきます。この直線を、点Pでの\textbf{接線}（せっせん）といい、Pを\textbf{接点}といいます。\bkidx{せっせん}{接線}

\begin{bkfigure}
\begin{tikzpicture}
\begin{axis}[bkplot, width=82mm, height=62mm, xmin=-0.3, xmax=2.4, ymin=-0.3, ymax=4.6, xtick={1,2}, ytick={1,2,3,4}, xlabel={$x$}, ylabel={$y$}]
  % 第6回：直線はx軸より上・枠の中で止め、名前は各直線の右の端に置く
  \addplot[bkcurve, domain=-0.3:2.1] {x^2};
  \addplot[bkthin, domain=0.72:2.2] {1+3*(x-1)};
  \addplot[bkthin, domain=0.64:2.4] {1+2.5*(x-1)};
  \addplot[bkthin, domain=0.6:2.35] {1+2.2*(x-1)};
  \addplot[bkcurve2, line width=1.4pt, domain=0.55:2.3] {1+2*(x-1)};
  \addplot[only marks, mark=*, mark size=1.8pt, bkInk] coordinates {(1,1) (2,4) (1.5,2.25) (1.2,1.44)};
  \node[anchor=north west, font=\sffamily\small] at (axis cs:1,1) {P};
  \node[anchor=south east, font=\sffamily\scriptsize, inner sep=1.5pt] at (axis cs:2,4) {Q};
  \node[anchor=south, font=\sffamily\scriptsize, text=bkGray, inner sep=1.5pt] at (axis cs:2.2,4.6) {傾き3};
  \node[anchor=west, font=\sffamily\scriptsize, text=bkGray, xshift=1mm] at (axis cs:2.4,4.5) {傾き2.5};
  \node[anchor=west, font=\sffamily\scriptsize, text=bkGray, xshift=1mm] at (axis cs:2.35,3.97) {傾き2.2};
  \node[anchor=west, font=\sffamily\scriptsize\bfseries, text=bkOrange!80!black, xshift=1mm] at (axis cs:2.3,3.6) {接線（傾き2）};
\end{axis}
\end{tikzpicture}
\bkfigcaption{$y=x^2$ の点P（1，1）。Qを（2，4）→（1.5，2.25）→（1.2，1.44）と近づけると、直線PQは接線（オレンジ）に近づく。}
\end{bkfigure}

\noindent\textbf{もう1つの見方：拡大すると、まっすぐに見える}\par
曲線を、点Pのまわりで虫めがねのように拡大していきます。下の3つの図は、どれも同じ大きさのわくに、Pのまわりをかいたものです。左から、横はば2、0.4、0.08の部分です。

\begin{bkfigure}
\begin{tikzpicture}
  % 3つのわくは同じ大きさ（38 mm 四方）。中心は P(1,1)。青い曲線 y=x^2、オレンジの点線は接線 y=2x-1
  \foreach \w/\xs/\lab in {1/0/横はば2, 0.2/46/横はば0.4, 0.04/92/横はば0.08} {
  \pgfmathsetmacro\xa{1-\w}\pgfmathsetmacro\xb{1+\w}\pgfmathsetmacro\ya{1-2*\w}\pgfmathsetmacro\yb{1+2*\w}
  \begin{axis}[at={(\xs mm,0)}, scale only axis, width=38mm, height=38mm, xmin=\xa, xmax=\xb, ymin=\ya, ymax=\yb,
    xtick=\empty, ytick=\empty, axis line style={bkGray, line width=0.6pt}, clip=true]
    \addplot[bkcurve, domain=\xa:\xb, samples=60] {x^2};
    \addplot[bkOrange, line width=0.9pt, dash pattern=on 3pt off 2pt, domain=\xa:\xb] {2*x-1};
    \addplot[only marks, mark=*, mark size=1.5pt, bkInk] coordinates {(1,1)};
  \end{axis}
  \node[font=\sffamily\scriptsize, anchor=north] at ({\xs mm+19mm},-2mm) {\lab};}
\end{tikzpicture}
\bkfigcaption{青が $y=x^2$、オレンジの点線が接線。拡大するほど、曲線は接線とほとんど重なる。}
\end{bkfigure}

拡大するほど、曲線はまっすぐに見え、接線とほとんど重なっていきます。つまり接線は、「その点のまわりをとても大きく拡大したとき、曲線が見せるまっすぐな線」です。だから、接線の傾きを、その点での「坂の急さ」と決めればよいのです。

\bksubsection{3}{接線の傾き：微分係数}

\leavevmode\label{bk:5-3:bibunkeisu}直線PQの傾きを、式で書きます。関数を $y=f(x)$ とし、Pの $x$ 座標を $a$、Qの $x$ 座標を $a+h$ とします。$a$ は「Pの場所」を表す決まった数で、$h$ は「PからQまでの横のはば」です。

\noindent\begin{minipage}[c]{0.5\linewidth}
\begin{tikzpicture}
\begin{axis}[bkplot, width=62mm, height=54mm, xmin=0, xmax=2.2, ymin=0, ymax=3.2, xtick={0.8,1.6}, xticklabels={$a$, $a+h$},
  ytick={0.64,2.56}, yticklabels={$f(a)$, $f(a+h)$}, xlabel={$x$}, ylabel={$y$}]
  \addplot[bkcurve, domain=0:1.75] {x^2};
  \draw[bkGray, line width=0.6pt, dash pattern=on 2pt off 1.5pt] (axis cs:0.8,0) -- (axis cs:0.8,0.64) -- (axis cs:0,0.64);
  \draw[bkGray, line width=0.6pt, dash pattern=on 2pt off 1.5pt] (axis cs:1.6,0) -- (axis cs:1.6,2.56) -- (axis cs:0,2.56);
  \addplot[bkcurve2, line width=1pt, domain=0.6:1.85] {0.64+2.4*(x-0.8)};
  \draw[bkInk, line width=0.7pt] (axis cs:0.8,0.64) -- (axis cs:1.6,0.64) -- (axis cs:1.6,2.56);
  \addplot[only marks, mark=*, mark size=1.7pt, bkInk] coordinates {(0.8,0.64) (1.6,2.56)};
  \node[anchor=south east, font=\sffamily\small] at (axis cs:0.8,0.64) {P};
  \node[anchor=east, font=\sffamily\small, xshift=-1.5mm] at (axis cs:1.6,2.56) {Q};
  \node[anchor=north, font=\sffamily\footnotesize] at (axis cs:1.2,0.6) {$h$};
\end{axis}
\end{tikzpicture}
\end{minipage}\hfill
\begin{minipage}[c]{0.47\linewidth}\raggedright
Pの高さは $f(a)$、Qの高さは $f(a+h)$。\par
横の変化量は $(a+h)-a=h$。\par
たての変化量は $f(a+h)-f(a)$（Qの右の、たての黒い線）。
\end{minipage}

\[ \text{PQの傾き}=\frac{\text{たての変化量}}{\text{横の変化量}}=\frac{f(a+h)-f(a)}{(a+h)-a}=\frac{f(a+h)-f(a)}{h} \]
これを\textbf{平均変化率}といいます。「$x$ が $a$ から $a+h$ まで変わるあいだに、$y$ が平均して $x$ の1あたりどれだけ変わったか」を表すからです。5-1の平均の速さは、$f$ が位置、$x$ が時間のときの平均変化率です。

Qを P に近づけることは、横のはば $h$ を0に近づけることです。だから、接線の傾きは
\[ f'(a)=\lim_{h\to0}\frac{f(a+h)-f(a)}{h} \]
です。これを、$x=a$ での\textbf{微分係数}（びぶんけいすう）といい、$f'(a)$（エフ・ダッシュ・エー）と書きます。$f$ の右上の「$'$」（ダッシュ）が「傾き」の印です。\bkidx{びぶんけいすう}{微分係数}\bkidx{へいきんへんかりつ}{平均変化率}

\noindent\textbf{やってみようの答え}\par
$f(x)=x^2$、$a=1$ とします。1行ずつ計算します。
\begin{align*}
f(1)&=1^2=1\\
f(1+h)&=(1+h)^2=1+2h+h^2 &&\text{（5-2 の正方形の図）}\\
f(1+h)-f(1)&=(1+2h+h^2)-1=2h+h^2\\
f'(1)&=\lim_{h\to0}\frac{2h+h^2}{h}=\lim_{h\to0}(2+h)=2
\end{align*}
定規で引いた直線の傾きは、ちょうど2です。図の傾き3（$h=1$）、2.5（$h=0.5$）、2.2（$h=0.2$）は、$2+h$ に $h$ を入れた値になっています。

\bksubsection{4}{どの点でも計算する：導関数}

\leavevmode\label{bk:5-3:doukansuu}$f(x)=x^2$ の微分係数を、いろいろな点で求めてみましょう。$a=1$ では2でした。同じ計算をすると、次のようになります（$a=3$ は例題2、$a=-1$ はじぶんでとこう）。

\begin{bkcenter}
\begin{bktable}{colspec={Q[l,m,wd=30mm] *{5}{Q[c,m,wd=13mm]}}, rowsep=1.5pt}
点の $x$ 座標 $a$ & $-1$ & 0 & 1 & 2 & 3 \\
微分係数 $f'(a)$ & $-2$ & 0 & 2 & 4 & 6 \\
\end{bktable}
\bktabcaption{$f(x)=x^2$ の、いろいろな点での接線の傾き}
\end{bkcenter}

どの点でも、$f'(a)$ は $a$ の2倍になっていそうです。1つずつ計算するかわりに、$a$ のところを $x$ のままにして計算すれば、どの点の接線の傾きも、一度に求められます。
\begin{align*}
f'(x)&=\lim_{h\to0}\frac{(x+h)^2-x^2}{h}=\lim_{h\to0}\frac{x^2+2xh+h^2-x^2}{h}\\
&=\lim_{h\to0}\frac{2xh+h^2}{h}=\lim_{h\to0}(2x+h)=2x
\end{align*}
表で見た「2倍」が、式で確かめられました。この $f'(x)$ を、$f(x)$ の\textbf{導関数}（どうかんすう）といいます。「$x$ を入れると、その点の接線の傾きが出てくる関数」です。導関数を求めることを、「$f(x)$ を\textbf{微分する}」といいます。$f(x)=x^2$ を微分すると $2x$ です。$x=1$ なら傾き2、$x=3$ なら傾き6、$x=-1$ なら傾き $-2$（右下がり）。\bkidx{どうかんすう}{導関数}\bkidx{びぶん}{微分}

\begin{bkfigure}
\begin{tikzpicture}
\begin{axis}[bkplot, name=top, width=62mm, height=46mm, xmin=-2.2, xmax=2.4, ymin=-0.5, ymax=4.6, xtick={-2,-1,1,2}, ytick={1,2,3,4}, xlabel={$x$}, ylabel={$y=x^2$}]
  \addplot[bkcurve, domain=-2.1:2.1] {x^2};
  % 第6回：接線は枠（ymax 4.6）の中で止める
  \foreach \a/\l/\r in {-1.5/-2/-1, 0/-0.5/0.5, 1/0.5/1.5, 2/1.6/2.15} {\edef\tmp{\noexpand\addplot[bkcurve2, line width=0.9pt, domain=\l:\r] {\a*\a+2*\a*(x-\a)};}\tmp}
\end{axis}
\begin{axis}[bkplot, at={($(top.east)+(12mm,0)$)}, anchor=west, width=62mm, height=46mm, xmin=-2.2, xmax=2.4, ymin=-4.4, ymax=4.6, xtick={-2,-1,1,2}, ytick={-4,-2,2,4}, xlabel={$x$}, ylabel={$y'=2x$}]
  \addplot[bkcurve2, domain=-2.1:2.1] {2*x};
  \addplot[only marks, mark=*, mark size=1.6pt, bkInk] coordinates {(-1.5,-3) (0,0) (1,2) (2,4)};
\end{axis}
\end{tikzpicture}
\bkfigcaption{左：$y=x^2$ と、いろいろな点の接線。右：傾きを点で打つと、直線 $y=2x$ にならぶ。}
\end{bkfigure}

導関数には、いくつかの書き方があります。

\begin{bkcenter}
\begin{bktable}{colspec={Q[c,m,wd=14mm] Q[l,m,wd=56mm] Q[l,m,wd=52mm]}, rowsep=1.5pt, cells={font=\small}, row{1}={font=\sffamily\bfseries\small}}
書き方 & 読み方 & よく使う場面 \\
$f'(x)$ & エフ・ダッシュ・エックス & 関数 $f$ を微分したとき \\
$y'$ & ワイ・ダッシュ & $y=\cdots$ の式を微分したとき \\
$\dfrac{\dd y}{\dd x}$ & ディーワイ・ディーエックス & $y$ を $x$ で微分したとき \\
$\dfrac{\dd x}{\dd t}$ & ディーエックス・ディーティー & 位置 $x$ を時間 $t$ で微分（速度） \\
\end{bktable}
\bktabcaption{導関数の書き方。どれも同じものを表す。}
\end{bkcenter}

$\dfrac{\dd y}{\dd x}$ は、平均変化率 $\dfrac{\Delta y}{\Delta x}$ の $\Delta$ を $\dd$ にかえた形です。「$\Delta x$ をかぎりなく小さくしたときの $\dfrac{\Delta y}{\Delta x}$」という意味を、形で表しています。分数のように見えますが、$\dd y$ を $\dd x$ でわった数ではなく、ひとまとまりの記号です。落ちるボール $x=5t^2$ なら、$\dfrac{\dd x}{\dd t}=10t$ が瞬間の速さです（5-2 の $v=10t$）。

\bksubsection{5}{接線の式}

\leavevmode\label{bk:5-3:shiki}接線は、点 $(a,\ f(a))$ を通り、傾きが $f'(a)$ の直線です。直線の式は、通る1点と傾きがわかれば作れます（3-5）。なぜそうなるのか、確かめておきます。

\noindent\begin{minipage}[c]{0.5\linewidth}
\begin{tikzpicture}
\begin{axis}[bkplot, width=62mm, height=52mm, xmin=0, xmax=4, ymin=-1.5, ymax=6.5, xtick={1,3}, xticklabels={$x_1$, $x$},
  ytick={1,5}, yticklabels={$y_1$, $y$}, xlabel={}, ylabel={}]
  \addplot[bkcurve2, line width=1.1pt, domain=0.3:3.6] {1+2*(x-1)};
  \draw[bkInk, line width=0.7pt] (axis cs:1,1) -- (axis cs:3,1) -- (axis cs:3,5);
  \addplot[only marks, mark=*, mark size=1.7pt, bkInk] coordinates {(1,1) (3,5)};
  \node[anchor=south east, font=\sffamily\footnotesize] at (axis cs:1,1) {A};
  \node[anchor=south east, font=\sffamily\footnotesize] at (axis cs:3,5) {B};
  \node[anchor=north, font=\sffamily\footnotesize] at (axis cs:2,0.9) {$x-x_1$};
  \node[anchor=west, font=\sffamily\footnotesize] at (axis cs:3.05,3) {$y-y_1$};
\end{axis}
\end{tikzpicture}
\end{minipage}\hfill
\begin{minipage}[c]{0.47\linewidth}\raggedright
傾き $m$ の直線が点A $(x_1,\ y_1)$ を通るとします。直線の上のほかの点を B $(x,\ y)$ とすると、AからBまでの傾きも $m$ です。
\[ \frac{y-y_1}{x-x_1}=m \]
両辺に $x-x_1$ をかけて（2-14）、
\[ y-y_1=m(x-x_1) \]
\end{minipage}

\medskip
数で確かめましょう。点 $(1,\ 1)$ を通り、傾き2の直線の上には、点 $(3,\ 5)$ があります（右へ2進むと、上へ $2\times2=4$ 上がる）。$\dfrac{5-1}{3-1}=\dfrac{4}{2}=2$ で、傾きは2です。

接線の場合は、$x_1$ が $a$、$y_1$ が $f(a)$、$m$ が $f'(a)$ なので、
\[ y-f(a)=f'(a)(x-a) \]
$y=x^2$ の点（1，1）での接線は、$f'(1)=2$ なので $y-1=2(x-1)$。右辺を展開して $y-1=2x-2$、両辺に1をたして $y=2x-1$ です。$x=1$ を入れると $y=1$ で、接点を通っています。

\begin{taisetsu}
微分係数（接線の傾き）\quad $f'(a)=\displaystyle\lim_{h\to0}\frac{f(a+h)-f(a)}{h}$\\
導関数 $f'(x)$ を求めることを「微分する」という。$(x^2)'=2x$\\
接線の式\quad $y-f(a)=f'(a)(x-a)$
\end{taisetsu}

\begin{kotoba}
\small
\bkword{極限}{かぎりなく近づけたときに、式の値が近づいていく先の値。}
\bkword{接線}{曲線の1点に、ぴったりふれる直線。その点のまわりを拡大すると、曲線と重なって見える。}
\bkword{微分係数}{1つの点での接線の傾き。$f'(a)$。}
\bkword{導関数}{$x$ を入れると、その点の接線の傾きが出てくる関数。$f'(x)$。}
\end{kotoba}

%% ---------------- 例題 ----------------
\bksection{reidai}{例題でたしかめよう}

\begin{reidai}[極限を求める]
$\displaystyle\lim_{h\to0}\frac{8h+2h^2}{h}$ を求めましょう。\par
\kangaekata $h=0$ を入れると $\dfrac{0}{0}$ になるので、$h$ を0に近づける前に、分子の1つずつを $h$ でわります。\par
\shiki $\dfrac{8h+2h^2}{h}=\dfrac{8h}{h}+\dfrac{2h^2}{h}=8+2h$\par
\hspace*{2.6\zw}$h$ を0に近づけると、$2h$ は0に近づくので、$8+2h$ は8に近づく。\par
\kotae 8
\end{reidai}

\begin{jibun}
$\displaystyle\lim_{h\to0}\frac{12h-3h^2}{h}$ を求めましょう。
\bkjans{$\lim(12-3h)=12$}
\end{jibun}

\begin{reidai}[微分係数を定義から求める]
$f(x)=x^2$ の、$x=3$ での微分係数 $f'(3)$ を求めましょう。\par
\kangaekata $f'(a)$ の式に $a=3$ を入れます。$f(3+h)$ は、$x$ に $3+h$ をかっこごと入れた値です。\par
\shiki $f(3)=3^2=9$\par
\hspace*{2.6\zw}$f(3+h)=(3+h)^2=9+6h+h^2$\par
\hspace*{2.6\zw}$f'(3)=\displaystyle\lim_{h\to0}\frac{(9+6h+h^2)-9}{h}=\lim_{h\to0}\frac{6h+h^2}{h}$\par
\hspace*{2.6\zw}$\phantom{f'(3)}=\displaystyle\lim_{h\to0}(6+h)=6$\par
\hspace*{2.6\zw}確かめ：$h=0.1$ なら $6+0.1=6.1$、$h=0.01$ なら $6.01$ で、6に近づく。\par
\kotae $f'(3)=6$
\end{reidai}

\begin{jibun}
$f(x)=x^2$ の、次の点での微分係数を、定義から求めましょう。\quad（1）$x=-1$\quad（2）$x=2$
\bkjans{（1）\nobreak $(-1+h)^2=1-2h+h^2$ より $\lim(-2+h)=-2$\quad（2）\nobreak $(2+h)^2=4+4h+h^2$ より $\lim(4+h)=4$}
\end{jibun}

\begin{reidai}[導関数を求める]
$f(x)=3x^2$ の導関数を、定義から求めましょう。\par
\kangaekata $f'(x)=\displaystyle\lim_{h\to0}\frac{f(x+h)-f(x)}{h}$ の分子を、先に計算します。\par
\shiki $f(x+h)=3(x+h)^2=3(x^2+2xh+h^2)=3x^2+6xh+3h^2$\par
\hspace*{2.6\zw}$f(x+h)-f(x)=3x^2+6xh+3h^2-3x^2=6xh+3h^2$\par
\hspace*{2.6\zw}$f'(x)=\displaystyle\lim_{h\to0}\frac{6xh+3h^2}{h}=\lim_{h\to0}(6x+3h)=6x$\par
\kotae $f'(x)=6x$
\end{reidai}

\begin{jibun}
$f(x)=x^2+x$ の導関数を、定義から求めましょう。
\bkjans{分子は $2xh+h^2+h$。$\lim(2x+h+1)=2x+1$}
\end{jibun}

\begin{reidai}[接線の式]
$y=x^2$ のグラフについて、（1）$x=2$ の点での接線の式を求めましょう。（2）傾きが6になる接線の、接点と式を求めましょう。\par
\kangaekata $(x^2)'=2x$ を使います。\par
\hspace*{4.3\zw}（1）接点は $(2,\ 2^2)=(2,\ 4)$、傾きは $f'(2)=2\times2=4$。\par
\hspace*{4.3\zw}（2）接点を $x=a$ とすると、傾きは $f'(a)=2a$。これが6になる $a$ を求める。\par
\shiki（1）$y-4=4(x-2)$\par
\hspace*{4.3\zw}$y-4=4x-8$\par
\hspace*{4.3\zw}$y=4x-8+4=4x-4$\par
\hspace*{2.6\zw}（2）$2a=6$ より $a=3$。接点は $(3,\ 3^2)=(3,\ 9)$\par
\hspace*{4.3\zw}$y-9=6(x-3)$ より $y=6x-18+9=6x-9$\par
\kotae （1）$y=4x-4$\quad（2）接点 $(3,\ 9)$、$y=6x-9$
\end{reidai}

\begin{jibun}
$y=x^2$ のグラフについて、（1）$x=-1$ の点での接線の式と、（2）傾きが2になる接線の接点と式を求めましょう。
\bkjans{（1）\nobreak 接点 $(-1,\ 1)$、傾き $-2$。$y=-2x-1$\quad（2）\nobreak $2a=2$ より接点 $(1,\ 1)$。$y=2x-1$}
\end{jibun}

\begin{reidai}[導関数から接線を引く]
$f(x)=x^2-4x$ の導関数を定義から求め、$x=1$ の点での接線の式を求めましょう。\par
\kangaekata $-4x$ の $x$ にも、かっこごと $x+h$ を入れます。\par
\shiki $f(x+h)-f(x)=(x+h)^2-4(x+h)-(x^2-4x)$\par
\hspace*{2.6\zw}$\phantom{f(x+h)-f(x)}=x^2+2xh+h^2-4x-4h-x^2+4x$\par
\hspace*{2.6\zw}$\phantom{f(x+h)-f(x)}=2xh+h^2-4h$\par
\hspace*{2.6\zw}$f'(x)=\displaystyle\lim_{h\to0}\frac{2xh+h^2-4h}{h}=\lim_{h\to0}(2x+h-4)=2x-4$\par
\hspace*{2.6\zw}接点 $(1,\ f(1))=(1,\ 1-4)=(1,\ -3)$、傾き $f'(1)=2\times1-4=-2$\par
\hspace*{2.6\zw}$y-(-3)=-2(x-1)$。$y+3=-2x+2$ より $y=-2x-1$\par
\kotae $f'(x)=2x-4$、$y=-2x-1$
\end{reidai}

\begin{jibun}
$f(x)=x^2+2x$ の導関数を定義から求め、次の点での接線の式を求めましょう。\quad（1）$x=1$\quad（2）$x=-2$
\bkjans{$f'(x)=2x+2$。（1）\nobreak 接点 $(1,\ 3)$、傾き4。$y=4x-1$\quad（2）\nobreak 接点 $(-2,\ 0)$、傾き $-2$。$y=-2x-4$}
\end{jibun}

\begin{reidai}[$x$-$t$ グラフの接線]
落ちるボールの位置は $x=5t^2$（m）で、$\dfrac{\dd x}{\dd t}=10t$ です。$x$-$t$ グラフの $t=1$ の点での接線の式を求めましょう。その傾きは何を表していますか。\par
\kangaekata 横軸が $t$、たて軸が $x$ なので、接線の式も $t$ と $x$ で書きます。「$y-f(a)=f'(a)(x-a)$」の $y$ が $x$ に、$x$ が $t$ にかわります。\par
\shiki 接点は $(1,\ 5\times1^2)=(1,\ 5)$、傾きは $10\times1=10$\par
\hspace*{2.6\zw}$x-5=10(t-1)$。$x-5=10t-10$ より $x=10t-5$\par
\noindent\begin{minipage}[c]{0.5\linewidth}
\begin{tikzpicture}
\begin{axis}[bkplot, width=56mm, height=40mm, xmin=0, xmax=2.2, ymin=-6, ymax=22, xtick={0,1,2}, ytick={-5,0,5,10,15,20},
  xlabel={$t$（s）}, ylabel={$x$（m）}]
  \addplot[bkcurve, domain=0:2] {5*x^2};
  \addplot[bkcurve2, line width=1pt, domain=0.1:2.2] {10*x-5};
  \addplot[only marks, mark=*, mark size=1.7pt, bkInk] coordinates {(1,5)};
\end{axis}
\end{tikzpicture}
\end{minipage}\hfill
\begin{minipage}[c]{0.46\linewidth}\raggedright
青が $x=5t^2$、オレンジが接線 $x=10t-5$。接線は点 $(1,\ 5)$ でグラフにふれています。$t=0$ を入れると $x=-5$ なので、たて軸と $-5$ で交わります。
\end{minipage}\par
\kotae $x=10t-5$。傾き10は、1秒の瞬間の速さ\qty{10}{m/s}を表す。
\end{reidai}

\begin{jibun}
$x=5t^2$（m）の $x$-$t$ グラフの、次の点での接線の式を求めましょう。\quad（1）$t=3$\quad（2）$t=0.5$
\bkjans{（1）\nobreak 接点 $(3,\ 45)$、傾き30。$x=30t-45$\quad（2）\nobreak 接点 $(0.5,\ 1.25)$、傾き5。$x=5t-1.25$}
\end{jibun}

%% ---------------- まちがえやすいところ・物理工学 ----------------
\bksection{tsumazuki}{まちがえやすいところをチェック}

\begin{tsumazuki}{$f(a+h)$ を $f(a)+h$ としてしまう}
$f(x)=x^2$ のとき、$f(3+h)$ は「$x$ に $3+h$ を入れた値」です。
\bkng{$f(3+h)=f(3)+h=9+h$}\bkok{$f(3+h)=(3+h)^2=9+6h+h^2$}
かっこごと入れるのが代入です（2-13）。$h=1$ で確かめると、$f(4)=16$ ですが、$9+1=10$ で合いません。
\end{tsumazuki}

\begin{tsumazuki}{接線の式で、傾きと接点を取りちがえる}
$y=x^2$ の $x=3$ の点の接線は、接点 $(3,\ 9)$、傾き $f'(3)=6$ です。
\bkng{$y-6=9(x-3)$}\bkok{$y-9=6(x-3)$}
かっこの前の数が傾き、左の $y$ からひく数が接点の $y$ 座標です。
\end{tsumazuki}

\begin{tsumazuki}{接線は、グラフと1点だけで交わる？}
\vspace{7pt}% 1行目が横に図のある段なので、題とのあいだをほかの箱とそろえる
\noindent\begin{minipage}[c]{20\zw}
円の接線は1点だけで円にふれます。でも、曲線の接線は、接点からはなれた所で、もう一度グラフと交わることがあります（右図）。接線かどうかは、「接点の近くで、グラフにぴったりふれているか」で決まります。
\end{minipage}\hfill
\begin{minipage}[c]{0.36\linewidth}
\begin{tikzpicture}
\begin{axis}[width=44mm, height=30mm, axis lines=none, xmin=0, xmax=6.5, ymin=-1.4, ymax=1.5, clip=false]
  % 第6回：x=4 の接線（枠の中で止める）。x=1.35 で、もう一度グラフと交わる
  \addplot[bkcurve, domain=0:6.3, samples=80, line width=1.2pt] {sin(deg(x))};
  \addplot[bkcurve2, domain=0.8:4.9, line width=1pt] {-0.7568-0.6536*(x-4)};
  \addplot[only marks, mark=*, mark size=1.4pt, bkInk] coordinates {(4,-0.7568)};
  \addplot[only marks, mark=o, mark size=1.6pt, bkInk, line width=0.8pt] coordinates {(1.35,0.976)};
  \node[anchor=south west, font=\sffamily\scriptsize, inner sep=1pt] at (axis cs:4.1,-0.48) {接点};
  \node[anchor=south west, font=\sffamily\scriptsize, inner sep=1pt] at (axis cs:1.6,1.3) {交わる};
\end{axis}
\end{tikzpicture}
\end{minipage}
\end{tsumazuki}

\begin{butsuri}
\begin{itemize}[leftmargin=1.2\zw, itemsep=2pt, topsep=0pt, label={\color{bkPurple}●}]
\item \textbf{ハンマー投げ}\par
回っているハンマーから手をはなすと、ハンマーはその瞬間の円の接線の向きに飛んでいきます。
\item \textbf{カーブの設計}\par
道路や線路は、まっすぐな所から曲がった所へ、向きが急に変わらないようにつないであります。つなぎめで、まっすぐな部分と曲がった部分の接線の向きが同じになるように設計するのです。
\item \textbf{速度計}\par
$x$-$t$ グラフの接線の傾きが、その瞬間の速さです。スピードメーターは、接線の傾きを（とても短い時間の平均で）計算する機械だといえます。
\end{itemize}
\end{butsuri}

\begin{motto}[3]
\textbf{とがった所には、接線が引けない}\par
$y=|x|$（$x$ の大きさ。$|{-3}|=3$）のグラフは、原点でとがっています。原点の右では傾き1、左では傾き $-1$ なので、右から近づけた傾きと左から近づけた傾きが一致しません。拡大しても、とがったままで、まっすぐには見えません。このような点では接線が決まらず、「微分できない」といいます。
\end{motto}

\begin{column}{ニュートンとライプニッツ}
微分の考え方は、350年ほど前に、イギリスのニュートンとドイツのライプニッツが、それぞれ別々に作りあげました。ニュートンは、惑星や落ちる物の運動を調べるために、時間で微分したものを $\dot{x}$（エックス・ドット）と、文字の上に点を打って書きました。ライプニッツは、$\dfrac{\dd y}{\dd x}$ という書き方を考えました。いまでも、物理では時間の微分に $\dot{x}$ を、数学や工学では $\dfrac{\dd y}{\dd x}$ をよく使います（5-5）。\bkidx{にゅーとん}{ニュートン}\bkidx{らいぷにっつ}{ライプニッツ}
\end{column}

%% ---------------- 章末問題 ----------------
\bkendproblems
\bkproblemgroup{基本}{3}

\bkproblem 次の極限を求めましょう。\quad（1）$\displaystyle\lim_{h\to0}\frac{5h+h^2}{h}$\quad（2）$\displaystyle\lim_{h\to0}\frac{4h-2h^2}{h}$
\bkpans{（1）\nobreak 5\quad（2）\nobreak 4}

\bkproblem $f(x)=x^2$ の、次の点での微分係数を定義から求めましょう。\\（1）$x=4$\quad（2）$x=-3$
\bkpans{（1）\nobreak $\lim(8+h)=8$\quad（2）\nobreak $\lim(-6+h)=-6$}

\bkproblem 次の関数の導関数を、定義から求めましょう。\\（1）$f(x)=2x^2$\quad（2）$f(x)=x^2-3x$
\bkpans{（1）\nobreak $4x$\quad（2）\nobreak $2x-3$}

\bkproblem $y=x^2$ のグラフの、$x=3$ の点での接線の式を求めましょう。
\bkpans{$y=6x-9$}

\bkproblem $f(x)=x^2+4x$ の導関数を定義から求め、$x=-1$ の点での接線の式を求めましょう。
\bkpans{$f'(x)=2x+4$。接点 $(-1,\ -3)$、傾き2。$y=2x-1$}

\bkproblem 次の極限を求めましょう。\quad（1）$\displaystyle\lim_{h\to0}\frac{3h+h^2}{h}$\quad（2）$\displaystyle\lim_{h\to0}\frac{-7h+4h^2}{h}$
\bkpans{（1）\nobreak 3\quad（2）\nobreak $-7$}

\bkproblemgroup{チャレンジ}{3}

\bkproblem 落ちるボールの位置は $x=5t^2$（m）で、$\dfrac{\dd x}{\dd t}=10t$ です。$x$-$t$ グラフの $t=2$ の点での接線の式を求めましょう。また、その傾きは何を表していますか。
\bkpans{$x=20t-20$。傾き20は、2秒の瞬間の速さ\qty{20}{m/s}}

\bkproblem $y=x^2$ のグラフの接線のうち、傾きが $-4$ になるものがあります。接点の座標と、接線の式を求めましょう。
\bkpans{$2a=-4$ より $a=-2$。接点 $(-2,\ 4)$、$y=-4x-4$}

\bkproblem $f(x)=x^2-2x$ の導関数を定義から求め、$x=3$ の点での接線の式を求めましょう。
\bkpans{$f'(x)=2x-2$。接点 $(3,\ 3)$、傾き4。$y=4x-9$}

\bkproblem 位置が $x=t^2-2t$（m）の物体があります。$x$-$t$ グラフの $t=3$ の点での接線の式を求め、その傾きが何を表すか答えましょう。
\bkpans{$\dfrac{\dd x}{\dd t}=2t-2$。接点 $(3,\ 3)$、傾き4。$x=4t-9$。傾きは3秒の瞬間の速さ\qty{4}{m/s}}

\bkproblem $f(x)=2x^2$ の導関数を定義から求め、$x=1$ の点での接線の式を求めましょう。
\bkpans{$f'(x)=4x$。接点 $(1,\ 2)$、傾き4。$y=4x-2$}

\begin{furikaeri}
\begin{checklist}
\item $\lim$ の記号を読み、極限の値を求められる。\hfill{\footnotesize\sffamily\color{bkTheme}→\pageref{bk:5-3:lim}ページ}
\item 微分係数を、接線の傾きとして求められる。\hfill{\footnotesize\sffamily\color{bkTheme}→\pageref{bk:5-3:bibunkeisu}ページ}
\item 導関数を求め、接線の式を書ける。\hfill{\footnotesize\sffamily\color{bkTheme}→\pageref{bk:5-3:shiki}ページ}
\end{checklist}
\end{furikaeri}

\bkfukidashi{delta}{$x^2$ を微分すると $2x$。では、$x^3$ や $x^4$ はどうなるかな？　次の章で、表をながめてきまりを見つけよう。}
